% !TEX program = pdflatex
% Обзор литературы по теме ВКР
% Компилятор в Overleaf: pdfLaTeX (Menu -> Compiler -> pdfLaTeX)
\documentclass[12pt,a4paper]{article}

% ---------- Кириллица и язык ----------
\usepackage{cmap}                    % копируемый русский текст в PDF
\usepackage[T2A]{fontenc}
\usepackage[utf8]{inputenc}
\usepackage[english,russian]{babel}

% ---------- Оформление страницы ----------
\usepackage[left=30mm,right=15mm,top=20mm,bottom=20mm]{geometry}
\usepackage{indentfirst}             % абзацный отступ у первого абзаца
\usepackage{setspace}
\onehalfspacing
\setlength{\parindent}{1.25cm}

% ---------- Формулы, таблицы, ссылки ----------
\usepackage{amsmath,amssymb}
\usepackage{booktabs}
\usepackage{array}
\usepackage{tabularx}
\usepackage[hidelinks]{hyperref}

\newcommand{\E}{\mathbb{E}}
\newcommand{\pos}[1]{\left(#1\right)^{+}}
\DeclareMathOperator*{\argmin}{arg\,min}

\title{\textbf{Обзор литературы}\\[4pt]
\large по теме выпускной квалификационной работы\\[2pt]
\large <<Вероятностное прогнозирование спроса и оптимизация закупок\\
с использованием интеллектуального агента>>}
\author{Студент: \underline{\hspace{6cm}}\\[2pt]
Научный руководитель: О.~А.~Тихонова}
\date{\today}

\begin{document}
\maketitle

\begin{abstract}
В обзоре рассмотрены пять работ 2016--2026~гг., посвящённых связке
<<прогноз спроса $\to$ решение о закупке>>. Показано, как задача
о газетчике (newsvendor) связывает вероятностный прогноз с оптимальным
объёмом заказа, как машинное обучение позволяет учитывать признаки
(цена, промо, праздники) без жёстких предположений о распределении
спроса и какую роль в управлении запасами начинают играть агентные
системы на основе больших языковых моделей. По итогам обзора
сформулированы пробелы в литературе и место данной ВКР среди
существующих работ.
\end{abstract}

\tableofcontents
\newpage

%=====================================================================
\section{Введение}
%=====================================================================

Закупка товара делается раньше, чем становится известен спрос.
Если заказать мало, магазин теряет маржу на неудовлетворённом спросе;
если заказать много, деньги замораживаются в запасе, а скоропортящийся
товар приходится списывать. Поэтому для закупок важен не один
<<самый вероятный>> прогноз, а распределение возможного спроса:
именно его хвосты определяют риск дефицита и размер страхового
запаса~\cite{silver2016inventory,hyndman2021fpp3}.

Классический подход разделяет задачу на два шага: сначала оценить
распределение спроса, затем по нему выбрать объём заказа.
В последние годы появились три направления, которые меняют эту схему:
\begin{enumerate}
  \item \textbf{обучение с признаками} --- объём заказа или квантиль
        спроса строится как функция от наблюдаемых факторов
        (день недели, цена, промо)~\cite{ban2019bigdata,oroojlooy2016deep};
  \item \textbf{вероятностные и генеративные модели спроса},
        позволяющие получить полное условное распределение и оценить
        риск решения~\cite{gong2024generative,ngartera2026stochastic};
  \item \textbf{агентные системы}, в которых прогноз, расчёт заказа
        и коммуникация с человеком объединены в один
        контур~\cite{syed2025agentic}.
\end{enumerate}

\textbf{Цель обзора} --- систематизировать эти подходы и обосновать
выбор методов для практической части ВКР: построение вероятностного
прогноза спроса, расчёт оптимального заказа и прототип агента,
объясняющего рекомендацию закупщику.

%=====================================================================
\section{Базовые понятия и обозначения}
%=====================================================================

\subsection{Задача о газетчике (newsvendor)}

Пусть $D$ --- случайный спрос за период, $q$ --- объём заказа,
$c_u$ --- потери от нехватки одной единицы (упущенная маржа),
$c_o$ --- потери от излишка одной единицы. Ожидаемые затраты
\begin{equation}\label{eq:nv}
  C(q) = \E\!\left[\, c_u \pos{D - q} + c_o \pos{q - D} \,\right],
  \qquad \pos{a} = \max\{a, 0\}.
\end{equation}
Приравняв производную к нулю, получаем, что оптимальный заказ
равен квантилю распределения спроса уровня \emph{критического
отношения}:
\begin{equation}\label{eq:cr}
  q^{*} = F^{-1}\!\left(\frac{c_u}{c_u + c_o}\right),
\end{equation}
где $F$ --- функция распределения спроса. Например, при цене
продажи $p = 1000$~тг и закупочной цене $c = 300$~тг имеем
$c_u = p - c = 700$, $c_o = c = 300$ и $q^{*} = F^{-1}(0{,}7)$.

\subsection{Квантильная функция потерь}

Для оценки квантиля уровня $\tau$ используется квантильная
(pinball) функция потерь~\cite{koenker1978quantiles}:
\begin{equation}\label{eq:pinball}
  L_\tau(y, \hat q) =
  \begin{cases}
    \tau\,(y - \hat q), & y \ge \hat q,\\[2pt]
    (1-\tau)\,(\hat q - y), & y < \hat q.
  \end{cases}
\end{equation}
Минимум $\E\,L_\tau(D, \hat q)$ по $\hat q$ достигается в квантиле
$F^{-1}(\tau)$. Если положить $\tau = c_u/(c_u+c_o)$, то
функция~\eqref{eq:pinball} с точностью до множителя совпадает
с затратами~\eqref{eq:nv}. Это ключевое наблюдение: \emph{хороший
квантильный прогноз уже является оптимальным решением о закупке}.
Pinball loss --- строго корректное правило оценки, поэтому её
используют для сравнения вероятностных прогнозов~\cite{gneiting2007scoring,makridakis2022m5}.

\subsection{Политика периодического контроля $(R, S)$}

Для товаров, которые не портятся, используется политика
<<раз в $R$ дней дозаказать до уровня $S$>>~\cite{silver2016inventory}.
Заказ равен $S - (\text{остаток} + \text{заказанное в пути})$,
а $S$ выбирается как квантиль спроса за период защиты $L + R$,
где $L$ --- срок поставки. Классическая формула при нормальном
спросе
\begin{equation}\label{eq:ss}
  S = \mu\,(L+R) + z_\alpha\,\sigma\sqrt{L+R}
\end{equation}
недооценивает запас при сверхдисперсии и прерывистом спросе
(классификация Syntetos--Boylan~\cite{syntetos2005accuracy}),
поэтому вероятностный прогноз здесь даёт выигрыш.

%=====================================================================
\section{Анализ отобранных работ}
%=====================================================================

\subsection{Ban, Rudin (2019): <<The Big Data Newsvendor>>}

\textbf{Задача.} Как выбрать объём заказа, если кроме истории
спроса есть $p$ признаков, известных до момента заказа~\cite{ban2019bigdata}.

\textbf{Метод.} Авторы предлагают не оценивать распределение,
а сразу искать решение как функцию признаков $q(x) = w^{\top}x$
методом минимизации эмпирического риска (ERM):
\begin{equation}
  \hat w = \argmin_{w}\; \frac{1}{n}\sum_{i=1}^{n}
  \Bigl[\, b\,\pos{d_i - w^{\top}x_i} + h\,\pos{w^{\top}x_i - d_i} \Bigr]
  + \lambda \lVert w \rVert_2^2 ,
\end{equation}
что эквивалентно линейной квантильной регрессии уровня $b/(b+h)$.
Второй алгоритм (Kernel Optimization, KO) взвешивает прошлые
наблюдения по близости признаков и решается сортировкой.

\textbf{Результаты.} Доказано, что игнорирование признаков приводит
к несостоятельным решениям, и получены оценки ошибки вне выборки
в зависимости от размерности. На данных о загрузке приёмного
отделения больницы в Великобритании лучшие алгоритмы ERM и KO
снизили затраты вне выборки на 23\% и 24\% относительно практического
эталона, причём KO оказался на три порядка быстрее ERM.

\textbf{Значение для ВКР.} Это теоретическое обоснование того,
почему в модель спроса нужно включать день недели, промо, праздники
и цену, и почему квантильная регрессия напрямую решает задачу закупки.

\subsection{Oroojlooyjadid, Snyder, Takáč (2016): глубокое обучение для newsvendor}

\textbf{Задача.} Многопризнаковая задача о газетчике при неизвестном
и изменчивом распределении спроса~\cite{oroojlooy2016deep}.

\textbf{Метод.} Глубокая нейронная сеть обучается с функцией потерь,
совпадающей с затратами newsvendor~\eqref{eq:nv}; выход сети --- сразу
объём заказа. Сеть не знает <<правильных>> меток, а выучивает их
в процессе минимизации затрат, то есть прогноз и оптимизация
объединены в один шаг.

\textbf{Результаты.} На реальных данных о продажах продуктовой корзины
сеть превосходит методы, основанные на оценке распределения, а также
kNN, случайный лес и ядерную регрессию, особенно при высокой
волатильности спроса. Для расширения на политику $(r, Q)$ полученное
решение оказалось на 19{,}9\% дешевле точного алгоритма EIL
и на 5{,}0\% дешевле kNN. Авторы отмечают, что при спросе из одного
хорошо описанного распределения и большом объёме данных разные
методы дают близкие результаты.

\textbf{Значение для ВКР.} Работа показывает, что выбор сложной
модели оправдан не всегда: на синтетических данных с известным
законом спроса (отрицательное биномиальное распределение) разумно
сравнить простую параметрическую модель с градиентным бустингом.

\subsection{Gong, Liu, Zhang (2024): генеративные модели спроса}

\textbf{Задача.} Задача о газетчике, в которой спрос зависит
от цены $p$ и признаков $x$, а требуется (а) выбрать запас при заданной
цене или (б) совместно выбрать цену и запас~\cite{gong2024generative}.

\textbf{Метод.} Условная глубокая генеративная модель (cVAE или cGAN)
обучается отображению $G_\theta(x, p, \eta)$, где
$\eta \sim \mathcal{N}(0, I)$ --- шум. Сгенерированные сценарии спроса
используются для оценки ожидаемой прибыли
\begin{equation}
  \hat\pi(x,p,q) = \frac{1}{M}\sum_{m=1}^{M}
  \Pi\bigl(G_{\hat\theta}(x,p,\eta_m),\, p,\, q\bigr),
\end{equation}
после чего решение находится максимизацией $\hat\pi$.

\textbf{Результаты.} Доказаны состоятельность оценки прибыли
и сходимость решений к оптимальным. В симуляциях (включая сценарий
с текстовыми признаками) и на открытом наборе Kaggle Food Demand
метод дал более высокую прибыль, чем SAA, KO и нейросетевой ERM.
Преимущество подхода --- одна обученная модель подходит для любых
цен и признаков; недостаток --- необходимость настройки нейросети.

\textbf{Значение для ВКР.} Идея <<сценарии спроса $\to$ оценка
прибыли методом Монте-Карло>> прямо используется в практической части
для получения распределения спроса за период $L + R$.

\subsection{Ngartera и др. (2026): вероятностный прогноз и риск}

\textbf{Задача.} Устранить разрыв между моделями прогноза и
операционными решениями о запасах с учётом ограничений по бюджету
и вместимости~\cite{ngartera2026stochastic}.

\textbf{Метод.} Предложена схема Decision Intelligence Framework
из трёх блоков: квантильная регрессия на градиентном бустинге;
задача о газетчике с ограничениями; оценка хвостового риска мерой
CVaR\textsubscript{95}:
\begin{equation}
  \mathrm{CVaR}_{\alpha}(X) = \E\bigl[\,X \mid X \ge \mathrm{VaR}_{\alpha}(X)\bigr].
\end{equation}
Формально показано соответствие <<калиброванный квантиль = решение
newsvendor>> (теорема~1) и субаддитивность CVaR по портфелю товаров
(теорема~2).

\textbf{Результаты.} На наборе UCI Online Retail (2010--2011)
при отношении штрафа за дефицит к стоимости хранения 10 политика по
квантилю P90 по сравнению с медианной P50 снизила CVaR\textsubscript{95}
на 26{,}7\%, повысила уровень обслуживания с 44{,}4\% до 89{,}5\%
и снизила средние затраты на 48{,}7\%. Авторы подчёркивают,
что выигрыш зависит от соотношения затрат.

\textbf{Значение для ВКР.} Работа ближе всего к практической части:
та же цепочка LightGBM (квантили) $\to$ критическое отношение
$\to$ симуляция затрат. Из неё берутся метрики (уровень обслуживания,
затраты, хвостовой риск) и вывод о необходимости калибровки квантилей.

\subsection{Syed и др. (2025): агентный ИИ для пополнения запасов}

\textbf{Задача.} Автоматизировать мониторинг запасов, прогноз,
выбор поставщика и заказ в розничной сети с множеством категорий
товаров, включая скоропортящиеся~\cite{syed2025agentic}.

\textbf{Метод.} Многоагентная архитектура AAIPS: агенты мониторинга,
прогноза спроса, решения о заказе, выбора поставщика, переговоров,
поиска трендовых товаров и координации. Политики агентов обучаются
многоагентным обучением с подкреплением (PPO); переговоры
предлагается вести с помощью больших языковых моделей.

\textbf{Результаты.} Заявлено снижение доли дефицитов примерно на 30\%,
затрат на хранение примерно на 15\% и общих затрат примерно на 10\%
относительно эвристик. При этом авторы прямо указывают, что
окончательные численные значения будут добавлены после
тестирования, а описание прототипа и данных краткое.
Результаты следует рассматривать как предварительные.

\textbf{Значение для ВКР.} Работа задаёт архитектурную рамку агента,
но показывает и риск: если числа порождает сама модель, их трудно
проверить. Поэтому в ВКР принят принцип <<все числа считают
инструменты, языковая модель только выбирает инструмент и объясняет
результат>>, а окончательное решение остаётся за человеком.

%=====================================================================
\section{Сравнительный анализ}
%=====================================================================

\begin{table}[h]
\centering
\small
\caption{Сравнение рассмотренных работ}
\label{tab:compare}
\begin{tabularx}{\textwidth}{@{}>{\raggedright}p{2.3cm}>{\raggedright}X>{\raggedright}X>{\raggedright}p{2.6cm}>{\raggedright\arraybackslash}X@{}}
\toprule
Работа & Модель спроса & Решение о заказе & Данные & Главный результат \\
\midrule
Ban, Rudin, 2019 & не оценивается; линейная по признакам & ERM / KO (квантильная регрессия) & больница, Великобритания & $-23\%$ и $-24\%$ затрат \\
Oroojlooyjadid и др., 2016 & не оценивается & DNN с функцией потерь newsvendor & продуктовая корзина & лучше при высокой волатильности \\
Gong и др., 2024 & условная генеративная (cVAE/cGAN) & максимизация прибыли по сценариям & симуляции, Kaggle Food Demand & выше прибыль, сходимость доказана \\
Ngartera и др., 2026 & квантили (бустинг) & критическое отношение, ограничения, CVaR & UCI Online Retail & P90: $-48{,}7\%$ затрат, сервис 89{,}5\% \\
Syed и др., 2025 & ML / временные ряды & MARL (PPO), агенты & данные магазина, симулятор & $\approx -30\%$ дефицитов (предв.) \\
\bottomrule
\end{tabularx}
\end{table}

Из таблицы~\ref{tab:compare} видно, что работы делятся на две группы.
\emph{Сквозные} методы~\cite{ban2019bigdata,oroojlooy2016deep} сразу
выдают объём заказа и не требуют предположений о распределении, но
плохо объясняют решение и привязаны к одной структуре затрат.
\emph{Двухшаговые} методы~\cite{gong2024generative,ngartera2026stochastic}
сначала строят распределение спроса, а затем решают задачу оптимизации;
они гибче (одно распределение годится для разных $c_u, c_o$ и
горизонтов) и дают интервалы, которые можно показать человеку.
Агентный подход~\cite{syed2025agentic} пока подтверждён слабее всего.

%=====================================================================
\section{Выводы и место ВКР}
%=====================================================================

\begin{enumerate}
  \item Оптимальный заказ в задаче о газетчике --- квантиль
        распределения спроса уровня $c_u/(c_u+c_o)$, поэтому
        вероятностный прогноз полезнее точечного для закупок;
        это подтверждается всеми рассмотренными работами.
  \item Учёт признаков (календарь, промо, цена) обязателен:
        без них решения несостоятельны~\cite{ban2019bigdata}.
  \item Сложные нейросетевые модели выигрывают при высокой
        волатильности и сложных признаках, но при достаточном объёме
        данных из известного распределения простые модели дают сопоставимый
        результат~\cite{oroojlooy2016deep}.
  \item Качество квантилей нужно проверять отдельно: калибровка,
        покрытие интервалов, pinball loss~\cite{gneiting2007scoring,ngartera2026stochastic}.
  \item Агентные системы для закупок находятся на ранней стадии;
        их слабое место --- проверяемость чисел и объяснимость решений~\cite{syed2025agentic}.
\end{enumerate}

\textbf{Пробел}, на который направлена ВКР: в рассмотренных работах
вероятностный прогноз, политика заказа и агентный интерфейс изучаются
по отдельности. В практической части ВКР они объединяются в один
воспроизводимый конвейер: <<данные $\to$ вероятностный прогноз
(квантильный LightGBM и отрицательная биномиальная GLM) $\to$
оптимальный заказ (newsvendor и $(R,S)$) $\to$ агент, объясняющий
решение>>. На синтетических данных с известным законом спроса можно
проверить, насколько точно модели восстанавливают распределение, а
симуляция на тестовом периоде позволяет оценить экономию в тенге.

\bibliographystyle{ugost2008}
\bibliography{references}

\end{document}
